% Part: sets-functions-relations
% Chapter: infinite
% Section: dedekind-induction

\documentclass[../../../include/open-logic-section]{subfiles}

\begin{document}

\olfileid[tr]{sfr}{infinite}{induction}
\olsection[Aritmetik Tümevarım]{Dedekind Cebirleri ve Aritmetik Tümevarım}

Artık önemli olan şudur: bir Dedekind cebiri---hatta \emph{herhangi bir}
Dedekind cebiri---doğal sayıların yerine kullanılabilir. Bunun nedeni
aşağıdaki dolaysız sonuçtur:

\begin{thm}[Aritmetik tümevarım]\ollabel{thm:dedinfiniteinduction}
$N,s,o$ bir Dedekind cebiri oluştursun. O zaman her $X$ kümesi için:
	\begin{center}
		$o \in X$ ve $(\forall {n} \in N \cap X){s}({n}) \in X$ ise
		$N \subseteq X$.
	\end{center}
\end{thm}

\begin{proof}
Bir Dedekind cebirinin tanımına göre $N=\closureofunder{s}{o}$'dır.
$Y=N\cap X$ olsun. $o\in N$ ve varsayım gereği $o\in X$ olduğundan
$o\in Y$'dir. Ayrıca $n\in Y$ ise hem $n\in N$ hem de $n\in X$ olur.
$N$, $s$ altında kapalı olduğundan $s(n)\in N$; varsayım gereği de
$s(n)\in X$ olur. Böylece $s(n)\in Y$ ve $Y$, $s$ altında kapalıdır.
$N=\closureofunder{s}{o}$, $o$'yu içeren en küçük $s$ altında kapalı küme
olduğundan $N\subseteq Y=N\cap X\subseteq X$.
\end{proof}

Tümevarım doğal sayıların ayırt edici bir özelliği olduğuna göre meselenin
özü şudur. Herhangi bir Dedekind-sonsuz küme verildiğinde bir Dedekind cebiri
oluşturabilir ve bu cebiri doğal sayıların yerine kullanabiliriz.

Kuşkusuz \olref{thm:dedinfiniteinduction}, tümevarımı \emph{küme teorisi}
terimleriyle formüle eder. Fakat ilkeyi daha tanıdık gelebilecek terimlerle
kolayca ifade edebiliriz:

\begin{cor}\ollabel{natinductionschema}
$N,s,o$ bir Dedekind cebiri oluştursun. O zaman parametreler içerebilen her
$\phi(x)$ formülü için:
\begin{center}
	$\phi(o)$ ve $(\forall {n} \in N)(\phi(n)\lif
	\phi({s}({n})))$ ise $(\forall n \in N)\phi(n)$.
\end{center}
\end{cor}

\begin{proof}
$X=\Setabs{n\in N}{\phi(n)}$ olsun ve şimdi
\olref{thm:dedinfiniteinduction}'ı kullanın.
\end{proof}

Bu sonuçta bir formülün ``parametreler içermesinden'' söz ettik. Bunun
kabaca anlamı, herhangi $c_1,\ldots,c_k$ nesneleri için
$\phi(x,c_1,\ldots,c_k)$ ile çalışabilmemizdir. Daha kesin olarak, sonucu
``parametreler''den söz etmeden şöyle ifade edebiliriz. Serbest değişkenlerinin
tümü gösterilmiş herhangi bir $\phi(x,v_1,\ldots,v_k)$ formülü için:
	\begin{align*}
			\forall v_1 \ldots \forall v_k((&\phi(o, v_1,\ldots, v_k) \land {}\\
			&	(\forall x \in N)(\phi(x,v_1, \ldots, v_k) \lif \phi(s(x), v_1,\ldots, v_k))) \lif {}\\
			&\hspace{3em} (\forall x \in N)\phi(x, v_1,\ldots, v_k))
	\end{align*}
elde ederiz. ``Parametreler içermek''ten söz etmek, açıkça, ifadeleri çok
daha okunaklı kılabilir. (Bu aracı ilerideki Küme Teorisi kısmında epey sık
kullanacağız.)

Dedekind cebirlerine dönelim: herhangi bir Dedekind cebiri verildiğinde,
toplama, çarpma ve üs alma gibi olağan aritmetik fonksiyonları da
tanımlayabiliriz. Ancak bu önemsiz değildir ve \emph{özyinelemeli tanım}
tekniğini gerektirir. Bu tekniği çok daha sonra ve çok daha genel bir
bağlamda tanıtıp gerekçelendireceğiz. (Meraklı okurlar ilerideki Sıra Sayısı
Aritmetiği bölümünden sonra buraya dönmek ya da
\citealt[pp.~95--8]{Potter2004} gibi alternatif bir sunumu okumak
isteyebilir.) Bununla birlikte, $N,s,o$ bir Dedekind cebiri oluşturduğunda
sonunda aşağıdakileri koşul olarak koyabileceğiz:
\begin{align*}
	{a} + {o} &= {a} & & & {a} \times {o} &= {o} & & & {a}^{o} &= s(o)\\
{a} + {s}({b}) &= {s}({a}+{b}) &&& {a} \times {s}({b}) &= ({a}\times {b}) + {a}  & & & {a}^{{s}({b})} &= {a}^{b} \times {a}
\end{align*}
ve bunların umduğumuz gibi davrandığını gösterebileceğiz.

\end{document}
